\section{Scope \& Methodology}
\subsection{In-Scope Sites}
\begin{center}\begin{tabularx}{\linewidth}{@{}l X@{}}
\toprule\rowcolor{TableHead}\textcolor{white}{\bfseries Site} & \textcolor{white}{\bfseries Areas In Scope} \\
\midrule
\fillme{Main office} & \fillme{Lobby, server room, executive floor} \\
\bottomrule\end{tabularx}\end{center}

\subsection{Rules of Engagement}
\begin{itemize}
  \item Testing window: \fillme{dates and times of day}.
  \item Authorisation letter: carried by all testers; \fillme{emergency contact name/number}.
  \item Destructive entry (forcing locks/doors): \fillme{in / out of scope}.
  \item Social engineering of staff: \fillme{in / out of scope, and any limits}.
\end{itemize}

\subsection{Methodology}
\begin{enumerate}
  \item \textbf{Reconnaissance.} Open-source and on-site observation of entry points, badge
        systems, camera coverage, and staff routines.
  \item \textbf{Perimeter testing.} Attempts to bypass fencing, doors, and other physical barriers.
  \item \textbf{Access control testing.} Tailgating, badge cloning (RFID), and lock-bypass
        (lock-picking, shimming) attempts against controlled entry points.
  \item \textbf{Post-access testing.} Once inside, attempts to reach sensitive areas (server
        rooms, network closets) and to plant an authorised test device (e.g. a drop box) to prove
        network access.
\end{enumerate}
